% ============================================================================
% Section 4 -- design. SECOND REWRITE 2026-09-07.
%
% The first rewrite added an architecture figure but was still hard to read:
% one static component diagram, dense prose, nothing concrete. A reader could
% not picture the compaction, could not see why k slots hide the transfer, and
% could not follow a request through the system.
%
% This version leads with four figures that carry the explanation:
%   fig:arch      the data path (components + who talks to whom)
%   fig:compact   the compaction, WORKED WITH REAL BLOCK IDS
%   fig:pipeline  the per-layer timeline -- why k>1 hides the transfer
%   fig:lifecycle a request's states and what triggers each transition
% Prose is trimmed to what the figures do not already say.
% ============================================================================

% PROVENANCE: the numbers in this section are design constants ($k{=}4$,
% $L{=}64$) or engine timings taken during development, as in
% \S\ref{sec:impl}. Measured results all live in \S\ref{sec:eval}, which
% carries the % src: comments.
\section{Design}
\label{sec:design}

\subsection{Overview}
\label{sec:design:overview}

Each request has a \emph{tier}: a \textbf{device-tier} request behaves exactly as
in stock vLLM, while a \textbf{host-tier} request holds \textbf{no blocks in the
device pool at all}. What follows is how that is possible --- why the host tier
needs its own address space, what a step does, and how the cohort is sized.
Appendix~\ref{app:figures} collects the notation, a component diagram and the
request-lifecycle state machine.

\subsection{Where the two tiers run}
\label{sec:design:exec}

\textbf{There is no partition of the GPU.} The two tiers are not separate streams,
kernels or forward passes but \emph{rows of the same batch}; the only operation that
distinguishes them is attention, and only in \emph{where it reads keys and values
from} (Fig.~\ref{fig:exec} draws one layer). Three consequences follow. Everything
but attention is shared, so a host-tier request adds one row to GEMMs that are
memory-bound on the weights regardless --- on an 8B model until batch 142
(Appendix~\ref{app:params}). The two attention calls are sequential on the same
hardware, so they do not compete for SMs, and the second's cost scales with the
host-tier rows rather than the batch. And the first does no wasted work for those
rows: masking their \texttt{seqused\_k} to 1 makes them attend a single key and be
overwritten by the second, worth 22\% of wall time. So a host-tier request occupies a
row in a step that was happening anyway, paying its KV read out of interconnect
bandwidth idle while the step reads weights --- the sense in which it rides along.

\subsection{Why the host tier has its own addresses}
\label{sec:design:spaces}

The host mirror is \textbf{not} indexed by device block id, and the rest of the
design rests on that. Under vLLM's allocator, releasing a block nulls the owner's
block-table entry and returns the id to a free list from which it is promptly
reissued --- so a mirror keyed by device block id would both \emph{lose the read
mapping} and \emph{alias}, the reissued id now denoting another request's data. Each
\emph{(request, logical block)} therefore gets an id from a private host space, held
until the request finishes, and device ids may churn underneath. Growth needs no
device allocation either: a new logical block takes a fresh host block and its
tokens arrive through the per-layer write path.

The price is that the cohort's host ids are scattered while the rolling region is a
small dense array. A \emph{compaction} once per step bridges them, sorting the
cohort's blocks onto $[0,R)$ so every layer reuses one block table while the slot
contents rotate beneath it.

\subsection{A decode step, end to end}
\label{sec:design:step}

Two granularities matter, and separating them keeps the per-layer work small:
\textbf{addresses} are fixed for the whole step, while slot \textbf{contents} change
every layer. Once per step, compact. Per layer, issue the runs into rolling slot
$\ell \bmod k$ --- with $k>1$ issued ahead, so layer $\ell{+}1$'s transfer overlaps
layer $\ell$'s compute, which is the hiding window the design turns on. Per layer,
attend: device-tier rows read the paged pool with the ordinary block table,
host-tier rows read slot $\ell \bmod k$ with the compacted one, in one pass over
shared weights. Per layer, write: the new $K,V$ for host-tier rows go to the
authoritative host mirror \emph{and} into the current slot, so this step's own
attention sees them. Several transitions are driven by the \emph{worker} rather than
the scheduler, which is not incidental: the scheduler cannot observe a host-side
write, so anything conditioned on KV having arrived must wait for an acknowledgement
(Appendix~\ref{app:figures}).

\subsection{Admission}
\label{sec:design:admission}

Each host-tier request costs $S\btok$ bytes of interconnect per step, and
admission bounds their sum by the step time times the achieved bandwidth, less a
margin. Three scheduler behaviours must then agree for such a request to hold no
device blocks, and all three are necessary. \emph{Admission without allocation}:
the request enters with an empty block list, its KV having been delivered into
the host pool by the prefill transport (\S\ref{sec:background:pd}), contingent on
the worker's acknowledgement. \emph{No prefill allocation}: its prompt is marked
already computed, so nothing is allocated even transiently --- the batch is
bounded by \emph{peak} device residency, so allocating during prefill and
reclaiming afterwards still pays the peak. \emph{Queue bypass}: vLLM's waiting
queue stops at the first request it cannot allocate for, so a block-starved
request ahead would otherwise block an admission that needs no resources.

\subsection{Placement and migration are mechanisms, not contributions}
\label{sec:design:placement}

\S\ref{sec:model} fixes \emph{how many} bytes each tier carries without reference
to which requests supply them: in the fluid limit, freed device memory and spent
interconnect both scale with $S\btok$, so any assignment filling both budgets
realises the whole gain. Placement can matter only through what that limit omits
--- integer effects at long context, remaining lifetime, KV growth, migration
cost, recompute-preemption --- so we present it as a mechanism the manager needs
rather than as a result.

The intuitive signal is the wrong one. \textbf{Context length does not discriminate
to first order}: moving a request of context $S$ frees $S\btok$ bytes and costs
$S\btok$ bytes every step, so ranking by context is ranking arbitrarily. What
survives is \emph{remaining} generation, which sets how many steps a one-time seed
amortises over; we estimate it as \texttt{max\_tokens} minus tokens produced --- a
bound the scheduler already holds, loose in the safe direction --- and admit above a
threshold at the amortisation point. On a substrate whose budget admits two host-tier
requests no ranking separates from arrival order, and \S\ref{sec:eval:ablation} says
so.

\label{sec:design:migration}
Placement must also be revisable, for two situations a fixed assignment handles
badly: device blocks free when a resident request finishes, and a host-tier request
promoted into them sheds its link traffic for nothing; and when the resident cohort
outgrows the pool, demotion is worth doing whenever it costs less than the
recompute-preemption it avoids. Both moves carry hysteresis and a cooldown bounding
thrash. The gate cannot be the byte budget alone, which prices transfer bytes while
the host path also carries a fixed per-step cost (\S\ref{sec:model:omits}), so a
demotion is a \textbf{probe}: make the move, settle, compare. The metric must be
goodput rather than latency, because a demotion \emph{always} lengthens the step ---
scored on latency the probe would revert every time --- while its benefit is
capacity, which shows up as requests served. On a link too slow to hide the transfer
no freed capacity repays the step it costs, so the probe should revert; on a
coherent link it should keep. That is a prediction:
\S\ref{sec:eval:ablation} explains why a two-rider budget cannot test it.
